\section{Related Work}
\label{sec:related_work}

\textbf{Model Merging in Deep Learning.} 
Weight averaging and model merging have gained traction as simple, training-free approaches to combine multiple specialized neural networks without increasing computational or memory costs during inference. Mitchell~\cite{mitchell80} and early studies highlighted the importance of inductive bias, while modern architectures have enabled direct linear interpolation of weights. Wortsman et al.~\cite{wortsman22} proposed \emph{Model Soups}, averaging multiple fine-tuned models of the same architecture to improve out-of-distribution generalization. Ilharco et al.~\cite{ilharco22} introduced \emph{Task Arithmetic}, which uses weight subtraction to create task-specific vectors that can be added to or subtracted from a pre-trained base model to steer model behaviors. 

To resolve parameter interference during merging, advanced techniques like \emph{Fisher Merging}~\cite{matena2022merging} weight parameters using diagonal Fisher information, while \emph{RegMean}~\cite{jin2023dataless} solves a closed-form linear system based on activation covariances. \emph{TIES-Merging}~\cite{yadav2023ties} resolves sign conflicts and parameter interference, and \emph{DARE}~\cite{yu2024language} randomly drops and rescales delta parameters to minimize interference. Modulo permutation symmetries, methods such as \emph{Git Re-Basin}~\cite{ainsworth2023git} and \emph{ZipIt!}~\cite{stoica2023zipit} align the weights of different initializations before averaging. Additionally, evolutionary merging recipes~\cite{akiba2025evolutionary} find optimal layer-wise combinations of weights and routing. We refer the reader to Yang et al.~\cite{yang2026ModelMergingSurvey} for a comprehensive survey of model merging methodologies. Despite their success, these standard methods operate entirely in high precision (FP32/FP16) and do not account for the severe discretization errors introduced by quantization.

\textbf{Quantization-Aware Model Merging.}
To bridge the gap between model merging and on-device deployment constraints, recent works have explored merging directly under quantized regimes. \emph{Q-Merge}~\cite{qmerge} optimizes layer-wise merging coefficients using the Straight-Through Estimator (STE)~\cite{bengio13ste} to propagate gradients through hard rounding operations. Similarly, \emph{ZipMerge}~\cite{zipmerge} combines joint weight pruning and test-time coefficient tuning to achieve highly compressed, sparse multi-task models. \emph{RegCalMerge}~\cite{regcalmerge} introduces test-time calibration and spatial regularization to prevent overfitting during coefficient tuning. However, a major shortcoming of these methods is their reliance on a single, static quantization operator (usually symmetric per-channel) during optimization. The robustness audit conducted in \cite{qmergeaudit} highlighted that such single-schema optimization makes the model highly vulnerable to catastrophic cross-schema overfitting when deployed onto hardware running mismatched compilers. Our work is the first to directly address and resolve this cross-schema generalization gap.

\textbf{Post-Training Quantization and Test-Time Adaptation.}
Post-training quantization (PTQ) is a standard technique to compress neural networks post-training without requiring expensive retraining~\cite{jacob18quant, gholami21survey, nagel19datafree}. Advanced PTQ methods like \emph{AdaRound}~\cite{nagel20adaround} optimize weight rounding to minimize PTQ reconstruction loss, while block-reconstruction methods like \emph{BRECQ}~\cite{li2021brecq} minimize the local output error. For large models, modern PTQ techniques like \emph{SmoothQuant}~\cite{pmlr-v202-xiao23c}, \emph{AWQ}~\cite{lin2024awq}, \emph{GPTQ}~\cite{frantar2023gptq}, \emph{LLM.int8()}~\cite{dettmers2022llmint8}, and \emph{QLoRA}~\cite{dettmers2023qlora} adapt scales, utilize outlier preservation, or introduce double quantization to maintain high fidelity under low-bit regimes. 

At the same time, Test-Time Adaptation (TTA) frameworks like \emph{Tent}~\cite{wang20tta} and \emph{EATA}~\cite{niu2022eata} adapt model parameters to unseen test streams by minimizing unsupervised Shannon prediction entropy or fisher-regularized entropy (see Liang et al.~\cite{liang2023comprehensive} for a comprehensive survey). Guided by these paradigms, we leverage test-time entropy minimization over a tiny, unlabeled calibration stream to adapt merging coefficients. Crucially, instead of keeping the quantization grid static, we introduce stochastic co-optimization to ensure that the adapted coefficients are robust to any hardware-specific PTQ schema.
